% §19.3 Maps as Stochastic Matrices
%
% Monograph register: the clean half. Deterministic f: U -> V is a column-
% stochastic 0/1 matrix M_f with M_f |x> = |f(x)>; the permutation special case.
% A Bernoulli map is a column-stochastic Q with Q_{yx} = P(ftilde(x)=y), so Q|x>
% is the output distribution; serial composition of Bernoulli maps is the matrix
% product Q_g Q_f, and the composition theorem omega = 1-(1-w1)(1-w2) is its
% consequence (Cref thm:composition / rem:channel-composition). Model order =
% number of distinct columns. This IS the n-ary channel of sec:n-ary-channels,
% read as an operator (note the row-stochastic vs column-stochastic convention
% difference: ch9 uses row-stochastic Q acting on the right; here column-
% stochastic acting on the left, the transpose). Cross-refs
% sec:operator-computational-basis, def:computational-basis, sec:n-ary-channels,
% def:n-ary-channel, sec:composition-theorem, thm:composition,
% rem:channel-composition, def:model-order, sec:operator-sets-tensored.
% DEFINED: sec:operator-maps-stochastic, def:operator-map-matrix,
% def:operator-bernoulli-map, prop:operator-map-composition.

\section{Maps as Stochastic Matrices}
\label{sec:operator-maps-stochastic}

A function is the cleanest thing to put in operator form, because a function
moves an element to an element, and an element is a basis vector.

\subsection*{A deterministic function is a $0/1$ matrix}

\begin{definition}[Map matrix]
\label{def:operator-map-matrix}
For a function $f : \Universe \to V$ between finite sets, the \emph{map matrix}
$M_f$ is the $\Card{V} \times \Card{\Universe}$ matrix with a single $1$ in each
column, placed so that
\[
  M_f \ket{x} \;=\; \ket{f(x)}
  \qquad \text{for every } x \in \Universe.
\]
Column $x$ of $M_f$ is the one-hot vector $\ket{f(x)}$; entry $(y, x)$ is $1$
when $y = f(x)$ and $0$ otherwise.
\end{definition}

Each column sums to one (a function sends each input somewhere, and to exactly
one place), so $M_f$ is \emph{column-stochastic} with entries in $\{0, 1\}$.
Applying $M_f$ to a basis vector reads off the function value; applying it to a
probability vector $p$ pushes the distribution forward, $M_f p$ being the law of
$f(X)$ when $X \sim p$. When $\Universe = V$ and $f$ is a bijection, $M_f$ is a
permutation matrix, the operator face of the user's $f(x) = M\ket{x}$ reading of
a function: a relabeling of the basis. A non-injective $f$ collapses several
basis vectors onto one, and $M_f$ then has a repeated column and is singular.

\subsection*{A Bernoulli map is a column-stochastic matrix}

Allowing the columns to be distributions rather than point masses turns a
function into an approximate one.

\begin{definition}[Bernoulli map matrix]
\label{def:operator-bernoulli-map}
A \emph{Bernoulli map} $\tilde{f} : \Universe \to V$ is given by a
column-stochastic matrix $Q$ of shape $\Card{V} \times \Card{\Universe}$,
\[
  Q_{yx} \;=\; \Prob{\tilde{f}(x) = y},
  \qquad
  Q_{yx} \ge 0,
  \quad
  \sum_{y} Q_{yx} = 1,
\]
so that column $x$, namely $Q \ket{x}$, is the output distribution on $V$
produced by the input $x$. The \emph{model order} (\Cref{def:model-order}) is
the number of distinct columns of $Q$, the number of input blocks that produce
genuinely different output distributions.
\end{definition}

A deterministic $f$ is the special case in which every column is one-hot, so
\Cref{def:operator-map-matrix} sits inside \Cref{def:operator-bernoulli-map}.
The matrix here is the n-ary channel of \Cref{def:n-ary-channel}, read as an
operator. The two definitions are the same object under transposition: the
channel of \Cref{sec:n-ary-channels} is written row-stochastic, with input
blocks indexing rows and the matrix acting on row vectors from the right, the
natural convention when one reads the matrix as a confusion table; the operator
here is written column-stochastic, with inputs indexing columns and the matrix
acting on column vectors $\ket{x}$ from the left, the natural convention when
one reads the matrix as a map. The entries are the same conditional
probabilities $\Prob{\text{output} \mid \text{input}}$, and ``column-stochastic
operator'' and ``row-stochastic channel'' name one matrix in two layouts. The
binary case $\Card{\Universe} = \Card{V} = 2$ recovers the $2 \times 2$ channel
of Part~II, with $Q \ket{0}$ and $Q \ket{1}$ the two columns carrying
$(1 - \fprate, \fprate)$ and $(\fnrate, 1 - \fnrate)$.

\subsection*{Composition is the matrix product}

The reason to write a map as a matrix is that composing maps becomes
multiplying matrices.

\begin{proposition}[Serial composition of Bernoulli maps]
\label{prop:operator-map-composition}
Let $\tilde{f} : \Universe \to V$ and $\tilde{g} : V \to W$ be Bernoulli maps
with column-stochastic matrices $Q_f$ and $Q_g$, and suppose the noise of
$\tilde{g}$ is independent of the noise of $\tilde{f}$. Then the composite
$\tilde{g} \circ \tilde{f} : \Universe \to W$ is a Bernoulli map with
column-stochastic matrix
\[
  Q_{g \circ f} \;=\; Q_g\, Q_f,
  \qquad
  (Q_g Q_f)_{wx} \;=\; \sum_{v \in V} (Q_g)_{wv}\, (Q_f)_{vx}.
\]
\end{proposition}

\begin{proof}
The composite output law at input $x$ is $\Prob{\tilde{g}(\tilde{f}(x)) = w}$.
Conditioning on the intermediate value $v = \tilde{f}(x)$ and using the
independence of the two noise sources,
\[
  \Prob{\tilde{g}(\tilde{f}(x)) = w}
  = \sum_{v \in V} \Prob{\tilde{g}(v) = w}\, \Prob{\tilde{f}(x) = v}
  = \sum_{v \in V} (Q_g)_{wv} (Q_f)_{vx},
\]
the $(w, x)$ entry of $Q_g Q_f$. Each column sums to one because $Q_f$'s column
$x$ is a distribution and $Q_g$ maps distributions to distributions
($\sum_w \sum_v (Q_g)_{wv}(Q_f)_{vx} = \sum_v (Q_f)_{vx} = 1$), so $Q_g Q_f$ is
column-stochastic and is therefore again a Bernoulli map.
\end{proof}

This is the operator face of serial channel composition
(\Cref{rem:channel-composition}), transposed to the column-stochastic layout
(the row-stochastic statement there is $Q_f Q_g$, the same product read in the
other order). The composition theorem of \Cref{thm:composition} is its
specialization to the one-sided error direction. Take two one-sided stages on a
shared codomain, each correct with probability $1 - \fnrate_i$ on a true input
and never producing a spurious accept. The relevant entry of the product
$Q_g Q_f$ is the probability that both stages pass the truth through, which by
the product is $(1 - \fnrate_1)(1 - \fnrate_2)$, so the composite miss rate is
\[
  \fnrate_{g \circ f} \;=\; 1 - (1 - \fnrate_1)(1 - \fnrate_2),
\]
and for a $k$-stage chain $\fnrate_{\mathrm{total}} = 1 - \prod_i (1 -
\fnrate_i)$, exactly \Cref{eq:composition-correctness}. The set-theoretic
chapter proved this by induction on the chain; the operator view gets it as one
diagonal entry of a matrix product, and the inductive step is the associativity
of matrix multiplication. The two derivations agree because they are the same
derivation in two notations.

\subsection*{Deterministic maps as a sanity check}

The product specializes correctly at the deterministic end. If $f$ and $g$ are
ordinary functions, then $M_g M_f \ket{x} = M_g \ket{f(x)} = \ket{g(f(x))} =
\ket{(g \circ f)(x)}$, so $M_{g \circ f} = M_g M_f$: the product of two
permutation matrices is the permutation of the composite, and the product of two
$0/1$ map matrices is the $0/1$ matrix of the composed function. The
approximate composition law degenerates to function composition exactly when the
columns degenerate to point masses, which is the consistency one wants between
the noisy theory and the exact one it generalizes. The next section turns to the
case where this clean single-matrix picture is the wrong one.
